El presente capítulo tiene como propósito establecer las políticas, procedimientos y medidas de control que deben aplicarse al transporte, manipulación y almacenamiento de los Sistemas de Aeronaves No Tripuladas (UAS), sus componentes, repuestos y baterías utilizadas por SKYFILMS, garantizando su preservación física, seguridad operacional y cumplimiento normativo.

Estas disposiciones buscan evitar daños, deterioro o condiciones inseguras durante el movimiento, transporte o almacenamiento de las aeronaves y sus elementos asociados, tanto dentro de las instalaciones como durante operaciones externas.

\section{Alcance}

Este procedimiento aplica a:

\begin{itemize}
    \item Todas las aeronaves UAS registradas bajo el certificado de operación de SKYFILMS.

    \item Los componentes electrónicos, estructurales, ópticos y de control asociados.

    \item Las baterías de polímero de litio (LiPo) y demás elementos energéticos.

    \item Las actividades de transporte terrestre, almacenamiento temporal y conservación prolongada de equipos.
\end{itemize}

\section{Principios Generales}
El transporte y almacenamiento de aeronaves y sus componentes deben realizarse bajo criterios de seguridad operacional, orden y trazabilidad.\\

Toda actividad en esta materia deberá cumplir los siguientes principios:

\begin{enumerate}
    \item \textbf{Protección física y ambiental:} las aeronaves y baterías deben transportarse y almacenarse protegidas contra golpes, vibraciones, humedad o temperaturas extremas.

    \item \textbf{Control documental:} cada traslado o almacenamiento deberá estar registrado, garantizando la trazabilidad del equipo y su estado técnico.

    \item \textbf{Responsabilidad técnica:} solo personal autorizado por el área de mantenimiento podrá manipular o autorizar movimientos de equipos UAS.

    \item \textbf{Cumplimiento normativo:} todas las actividades deben observar lo dispuesto en el RAC 100, el RAC 175 y las normas ambientales y de seguridad aplicables.
\end{enumerate}

\section{Transporte de Aeronaves UAS}

El transporte de aeronaves UAS es una actividad que requiere cuidado y supervisión técnica. Antes, durante y después del traslado, deben cumplirse las siguientes disposiciones:

\begin{itemize}
    \item \textbf{Preparación previa:} El técnico responsable verificará que el UAS se encuentre apagado, sin baterías instaladas, con hélices retiradas y brazos plegados o asegurados. Los equipos ópticos o sensores desmontables deben retirarse y guardarse por separado.

    \item \textbf{Embalaje y protección:} El transporte se realizará en maletas rígidas, estuches acolchados o contenedores de seguridad, diseñados para absorber impactos y evitar movimientos bruscos.

    \item \textbf{Transporte terrestre y externo:} Durante el traslado, los equipos viajarán en vehículos con superficies estables, asegurados mediante correas o espuma protectora. Las condiciones ambientales del transporte deben garantizar ausencia de humedad y exposición directa al sol.

\end{itemize}

Estas medidas aseguran que cada aeronave conserve sus condiciones estructurales y electrónicas sin alteración durante los traslados.

\section{Transporte y Manipulación de Baterías}

Las baterías de polímero de litio (LiPo) son elementos de alta sensibilidad y requieren una manipulación controlada. El procedimiento de transporte debe contemplar las siguientes medidas:

\begin{itemize}
    \item \textbf{Condición de carga:} Toda batería debe transportarse con un nivel máximo del 30\% al 40\% de su capacidad, reduciendo el riesgo térmico.

    \item \textbf{Aislamiento:} Las baterías se guardarán en bolsas ignífugas o contenedores antiestáticos.

    \item \textbf{Condiciones ambientales:} No se permitirá el transporte en ambientes cerrados sin ventilación ni la exposición a fuentes de calor.

\end{itemize}

El cumplimiento de estas prácticas garantiza la seguridad del personal y la integridad de las baterías durante el transporte.

\section{Almacenamiento de Aeronaves UAS}
Las aeronaves UAS deben conservarse en un entorno controlado que asegure su integridad estructural y electrónica, minimizando los efectos del ambiente o la manipulación inadecuada. Para ello, se establecen las siguientes disposiciones:

\begin{enumerate}
    \item \textbf{Condiciones del área:} El área de almacenamiento debe ser seca, limpia, ventilada y libre de polvo.

    \item \textbf{Distribución y señalización:} Cada aeronave dispondrá de un espacio asignado con etiqueta de identificación que indique su modelo, número de serie y estado técnico (Operativa, En Mantenimiento, No Aeronavegable).

    \item \textbf{Protección física:} Las aeronaves se colocarán sobre superficies planas y estables, evitando apilamientos o presión sobre las estructuras.

    \item \textbf{Inspecciones periódicas:} El Técnico UAS verificará \textbf{mensualmente} el estado del área y las aeronaves almacenadas, registrando cualquier hallazgo o deterioro.

\end{enumerate}

Estas acciones aseguran la preservación de los equipos cuando no están en uso o se encuentran en espera de intervención técnica.


\section{Almacenamiento de Baterías}
Las baterías de litio son elementos críticos que requieren una gestión separada y condiciones controladas para su seguridad. El almacenamiento debe cumplir los siguientes requisitos:

\begin{itemize}
    \item \textbf{Área designada y aislada:} Las baterías se almacenarán en un área independiente, ventilada, con señalización visible de riesgo eléctrico e incendio.

    \item \textbf{Condiciones ambientales:} Temperatura entre 15°C y 25°C, humedad relativa inferior al 60\%, sin exposición directa al sol ni a fuentes de calor.

    \item \textbf{Nivel de carga:} Para almacenamiento prolongado, las baterías deben mantenerse con una carga parcial (~40\%), y someterse a revisiones mensuales de voltaje y estado físico.

    \item \textbf{Control de seguridad:} El área deberá contar con detectores de humo, extintores tipo D o E para caso de incendio o fuga.
\end{itemize}

El Técnico UAS será responsable del control y registro de las baterías almacenadas, garantizando que ninguna batería deteriorada permanezca en el área.

\section{Responsabilidades}

\begin{itemize}
    \item \textbf{Gerente General:} Garantiza la disponibilidad de recursos, condiciones y equipos necesarios para cumplir las disposiciones de transporte y almacenamiento, asegurando el cumplimiento de las normas aeronáuticas y de seguridad.

    \item \textbf{Tecnico UAS:} 
    
    \begin{itemize}
        \item Supervisa la aplicación de los procedimientos establecidos, aprueba los traslados fuera de la base técnica y revisa los registros asociados.

        \item Administra los registros, controla las inspecciones de almacenamiento, autoriza movimientos y verifica la seguridad de las condiciones ambientales.

        \item Preparan las aeronaves para transporte, ejecutan inspecciones previas y reportan cualquier daño o desviación detectada.
    \end{itemize}


\end{itemize}

El cumplimiento conjunto de estas responsabilidades garantiza la trazabilidad, seguridad y preservación del material técnico de SKYFILMS.

\section{Disposiciones Finales}

El transporte y almacenamiento adecuado de aeronaves, componentes y baterías constituye un requisito esencial para la \textbf{seguridad operacional y la aeronavegabilidad continuada} de la flota UAS de SKYFILMS.\\

El incumplimiento de los procedimientos establecidos en este capítulo será considerado una no conformidad técnica grave, susceptible de acciones correctivas y sanciones administrativas internas. El seguimiento constante y la correcta aplicación de estos lineamientos aseguran que cada equipo mantenga sus condiciones originales, garantizando operaciones seguras, sostenibles y conformes con los Reglamentos Aeronáuticos de Colombia.